\documentclass[11pt]{article}
\usepackage[margin=2cm,headheight=14pt]{geometry}
\usepackage{amsmath,amssymb}
\usepackage{fancyhdr}
\usepackage[numbers]{natbib}

\pagestyle{fancy}
\fancyhf{}
\fancyhead[L]{Astroestadística}
\fancyhead[R]{Semestre 2026-2}
\renewcommand{\headrulewidth}{0.4pt}

\setlength{\parindent}{0pt}
\setlength{\parskip}{0.5\baselineskip}

\begin{document}
\section*{Problemset 1: transformación inversa}

Los problemas están ordenados para construir el método paso a paso. En los
problemas 1--6 se trabajará con una misma densidad; después se aplicarán la
inversa generalizada y la inversión numérica. Muestre los cálculos y justifique
sus respuestas.

\begin{enumerate}
    \item Considere la función no normalizada
    \[
        \widetilde f(x)=\begin{cases}x+2,&0\leq x\leq2,\\0,&\text{en otro caso.}\end{cases}
    \]
    Calcule su constante de normalización $c=\int_{-\infty}^{\infty}\widetilde
    f(x)\,dx$ y escriba la densidad normalizada $f(x)=\widetilde f(x)/c$.

    \textbf{Solución.} La función es distinta de cero únicamente entre $0$ y
    $2$, por lo que
    \[
        c=\int_0^2(x+2)\,dx
        =\left[\frac{x^2}{2}+2x\right]_0^2=6.
    \]
    En consecuencia,
    \[
        f(x)=\begin{cases}
            \dfrac{x+2}{6},&0\leq x\leq2,\\
            0,&\text{en otro caso.}
        \end{cases}
    \]
    La normalización es necesaria antes de aplicar la transformación a una
    distribución objetivo; véase la discusión de transformaciones en
    \cite[Sec.~7.2]{press1992numerical}.

    \item A partir de la densidad normalizada del problema 1, derive la función
    de distribución acumulada $F(x)$ para todo $x\in\mathbb{R}$. Compruebe
    que $F(0)=0$ y $F(2)=1$.

    \textbf{Solución.} Integramos la densidad desde el extremo izquierdo del
    soporte. Incluyendo los tramos fuera del soporte,
    \[
        F(x)=\begin{cases}
            0,&x<0,\\[2pt]
            \displaystyle\int_0^x\frac{t+2}{6}\,dt
            =\dfrac{x^2+4x}{12},&0\leq x\leq2,\\[6pt]
            1,&x>2.
        \end{cases}
    \]
    Se verifica que $F(0)=0$ y $F(2)=(4+8)/12=1$. La relación entre
    transformaciones y distribuciones objetivo se presenta en
    \cite[Sec.~7.2]{press1992numerical}.

    \item Despeje $x$ de la ecuación $u=F(x)$ en el intervalo $0\leq x\leq2$.
    Escriba la función inversa $F^{-1}(u)$ para $0\leq u\leq1$ y verifique
    que sus valores quedan dentro del soporte de la distribución.

    \textbf{Solución.} En el soporte, $u=(x^2+4x)/12$. Reordenando y
    resolviendo la ecuación cuadrática,
    \[
        x^2+4x-12u=0,
        \qquad x=-2\pm\sqrt{4+12u}.
    \]
    Se toma la raíz positiva porque la otra es negativa en el intervalo
    considerado. Así,
    \[
        F^{-1}(u)=-2+\sqrt{4+12u},\qquad 0\leq u\leq1.
    \]
    Como esta expresión es creciente, sus extremos son $F^{-1}(0)=0$ y
    $F^{-1}(1)=2$; por tanto, siempre cae en el soporte. La transformación de
    uniformes para obtener distribuciones se trata en
    \cite[Sec.~7.2]{press1992numerical}.

    \item Use la transformación inversa para generar una observación cuando
    $U=0.75$. Dé el resultado exacto y una aproximación con tres cifras
    decimales. Explique qué operación se aplica al número uniforme.

    \textbf{Solución.} Se sustituye el uniforme en la inversa del problema 3:
    \[
        X=F^{-1}(0.75)=-2+\sqrt{4+12(0.75)}
        =\sqrt{13}-2\approx1.606.
    \]
    Es decir, se transforma el número uniforme evaluando la función inversa de
    la FDA. Este es el principio de transformación descrito en
    \cite[Sec.~7.2]{press1992numerical}.

    \item Usando la FDA obtenida en el problema 2, calcule
    $\Pr(0.5<X\leq1.5)$. Exprese la respuesta como fracción y como número
    decimal.

    \textbf{Solución.} Como $X$ es continua,
    \[
        \Pr(0.5<X\leq1.5)=F(1.5)-F(0.5)
        =\frac{(1.5)^2+4(1.5)}{12}
         -\frac{(0.5)^2+4(0.5)}{12}
        =\frac{33}{48}-\frac{9}{48}
        =\frac12=0.5.
    \]
    Se ha usado la FDA obtenida al integrar la densidad; la transformación
    asociada a la distribución objetivo se discute en
    \cite[Sec.~7.2]{press1992numerical}.

    \item Encuentre la mediana $m$ de esta distribución, es decir, el valor que
    satisface $F(m)=0.5$. Interprete el resultado en términos de probabilidad.

    \textbf{Solución.} Sustituimos $u=0.5$ en la inversa:
    \[
        m=F^{-1}(0.5)=-2+\sqrt{4+6}=\sqrt{10}-2
        \approx1.162.
    \]
    Por definición, $F(m)=0.5$: la mitad de la probabilidad se encuentra a
    cada lado de la mediana, es decir, $\Pr(X\leq m)=0.5$. La transformación
    inversa se aborda en \cite[Sec.~7.2]{press1992numerical}.

    \item Sea ahora $Y$ una variable discreta con
    \[
        \Pr(Y=-1)=0.2,\qquad \Pr(Y=2)=0.35,\qquad \Pr(Y=6)=0.45.
    \]
    Escriba su inversa generalizada $F_Y^{-1}(u)$ como una función por tramos.
    Luego determine el valor de $Y$ que se obtiene para $U=0.4$ y para $U=0.9$.

    \textbf{Solución.} Las probabilidades acumuladas en los puntos de soporte
    son $0.2$, $0.2+0.35=0.55$ y $1$. Usando la convención
    $F_Y^{-1}(u)=\inf\{y:F_Y(y)\geq u\}$,
    \[
        F_Y^{-1}(u)=\begin{cases}
            -1,&0<u\leq0.2,\\
            2,&0.2<u\leq0.55,\\
            6,&0.55<u\leq1.
        \end{cases}
    \]
    Por tanto, $F_Y^{-1}(0.4)=2$ y $F_Y^{-1}(0.9)=6$. Los intervalos se
    determinan directamente de la FDA discreta del enunciado; para el
    principio general de transformar uniformes véase
    \cite[Sec.~7.2]{press1992numerical}.

    \item Considere la densidad
    \[
        g(x)=\begin{cases}
            \dfrac{3}{4}(1+x^2),&0\leq x\leq1,\\
            0,&\text{en otro caso.}
        \end{cases}
    \]
    Su FDA en $[0,1]$ es $G(x)=\frac{3}{4}\bigl(x+\frac{x^3}{3}\bigr)$.
    Para generar una observación cuando $U=0.5$, plantee la ecuación que debe
    resolverse y use bisección para aproximar $G^{-1}(0.5)$ al centésimo.
    Indique por qué aquí conviene una inversión numérica.

    \textbf{Solución.} Igualamos la FDA a $0.5$ y simplificamos:
    \[
        \frac34\left(x+\frac{x^3}{3}\right)=\frac12
        \quad\Longleftrightarrow\quad x^3+3x-2=0.
    \]
    Sea $h(x)=x^3+3x-2$. Como $h(0)=-2$ y $h(1)=2$, hay una raíz en
    $[0,1]$. La bisección conserva en cada paso el subintervalo cuyos extremos
    dan signos opuestos. Tras nueve pasos se obtiene, por ejemplo,
    \[
        0.595703125\leq x_*\leq0.59765625,
    \]
    de modo que al centésimo $G^{-1}(0.5)\approx0.60$ (más precisamente,
    $x_*\approx0.59607$). La ecuación cúbica no tiene aquí una inversa
    práctica para el cálculo manual, por lo que la bisección ofrece una
    aproximación controlada. El método de bisección se explica en
    \cite[Sec.~9.1]{press1992numerical}.
\end{enumerate}

\bibliographystyle{plainnat}
\bibliography{referencias}

\end{document}